\section{可观测量}\label{sec:observables}

\gff\ 中可以考虑的最简单的量之一，是规范不变算符在有限流时间处的真空期望值。
\gff\ 的一个显著性质是：这些算符除常规 \qcd\ 的重正化以及所涉流动场自身的重
正化之外，不需要任何额外的重正化。对规范耦合与夸克质量，\msbar\ 重正化通过
在 \noeqn{eq:lags} 中作代换
\begin{equation}
  \begin{split}
    g\to g_0\equiv \left(\frac{\mu\mathrm{e}^{\gamma_\mathrm{E}/2}}{\sqrt{4\pi}}\right)^\ep
    Z_g(\alphas(\mu))\,g(\mu)\,,\qquad m_f\to m_{f,0}\equiv
    Z_m(\alphas(\mu))\,m_f(\mu)
    \label{eq:renorm}
  \end{split}
\end{equation}
实现，其中 $\mu$ 为重正化标度，$\alphas=g^2/(4\pi)$，
$\EulerGamma$ 为 Euler-Mascheroni 常数。在本文所需的微扰阶内，重正化常数为
\begin{equation}
  \begin{split}
    Z_g(\alphas) &=
    1 -
    \frac{\alphas}{4\pi}\frac{\beta_0}{2\ep} +
    \left(\frac{\alphas}{4\pi}\right)^2\left(\frac{3\beta_0^2}{8\ep^2}
    -\frac{\beta_1}{4\ep}\right) + \order{\alphas^3}\,, \\ Z_m(\alphas) &=
    1 -\frac{\alphas}{4\pi}\frac{\gamma_{m,0}}{2\ep}
    +\left(\frac{\alphas}{4\pi}\right)^2\left[\frac{1}{\ep^2}\left(
      \frac{\gamma_{m,0}^2}{8} + \frac{\beta_0\gamma_{m,0}}{4}\right)
      -\frac{\gamma_{m,1}}{4\ep}\right] + \order{\alphas^3}\,,
    \label{eq:rgconsts}
  \end{split}
\end{equation}
其中
\begin{equation}
\begin{aligned}
      \beta_0 &= \frac{11}{3} \CA - \frac{4}{3} \TF\,, &
      \beta_1 &= \frac{34}{3} \CA^2 - \left( 4 \CF + \frac{20}{3} \CA
      \right) \TF \,,\\
      \gamma_{m,0} &= 6 \CF\,, &
      \gamma_{m,1} &= \frac{97}{3} \CA \CF + 3 \CF^2
      - \frac{20}{3} \CF \TF \,.
    \label{eq:anombg}
\end{aligned}
\end{equation}
$\CF$ 与 $\CA$ 分别是规范群基本表示与伴随表示的二次 Casimir 本征值。此外，
$\TF=\TR \NF$，其中 $\NF$ 为夸克味的数目，$\TR$ 为基本表示的迹归一化。对
SU($\NC$)，有 $\CF=(\NC^2-1)/(2\NC)$、$\CA=\NC$、$\TR=1/2$。

流动规范场 $B^a_\mu(t,x)$ 不需要重正化，因此例如胶子作用量密度
\begin{align}
  E (t,x) \equiv \frac{1}{4} G_{\mu \nu}^a(t,x) G_{\mu \nu}^a(t,x)\,,
  \label{eq:def_E}
\end{align}
的矩阵元在仅对 $g$ 与 $m_f$ 作重正化之后即为有限。这使得不同正则化方案（例
如格点与微扰论）所得的结果可以直接比较。

相反，流动夸克场需要一个重正化因子 $Z_\chi^{1/2}(\alphas)$ 才能使格林函数有
限。在 \msbar\ 方案中它为\,\cite{Luscher:2010iy,Harlander:2018zpi}
\begin{equation}
  \begin{split}
        Z^{-1}_\chi(\alphas) &= 1
        -\frac{\alphas}{4\pi}\frac{\gamma_{\chi,0}}{2\ep}
        +\left(\frac{\alphas}{4\pi}\right)^2\left[\frac{1}{\ep^2}\left(
          \frac{\gamma_{\chi,0}^2}{8} +
          \frac{\beta_0\gamma_{\chi,0}}{4}\right)
          -\frac{\gamma_{\chi,1}}{4\ep}\right] + \order{\alphas^3}\,,
    %% \label{eq:}
  \end{split}
\end{equation}
其中
\begin{equation}
  \begin{split}
  \gamma_{\chi,0} &= 6\,\CF\,,\\
  \gamma_{\chi, 1} &=
  \left( \frac{223}{3} - 16 \ln 2 \right) \CA \CF - \left( 3 + 16
  \ln 2 \right) \CF^2 - \frac{44}{3} \CF \TF\,.
    %% \label{eq:}
  \end{split}
\end{equation}

于是，量
\begin{align}
  S (t,x) \equiv Z_\chi\,\sum_{f=1}^{\NF}\bar{\chi}_f(t,x) \chi_f(t,x)
  \label{eq:def_S}
\end{align}
获得一个反常量纲，这妨碍了不同正则化方案结果的直接比较。另一种做法是使用
``加环夸克场''（ringed quark fields）\,\cite{Makino:2014taa}，即改用
\begin{equation}
  \begin{split}
    \mathring{Z}_\chi(t,\mu) = -\frac{2\NC \NF}{(4\pi t)^2}\, \langle
    R(t)\rangle^{-1}\,,\quad\mbox{其中}\quad R(t,x) =
    \sum_{f=1}^{\NF}\bar{\chi}_f(t,x) \overleftrightarrow{\slashed{\mathcal{D}}}^\mathrm{F} \chi_f(t,x)
  \label{eq:Zring}
  \end{split}
\end{equation}
代替 $Z_\chi$ 来重正化夸克场，这里
\begin{equation}
  \overleftrightarrow{\mathcal{D}}^\mathrm{F} = \mathcal{D}^\mathrm{F} - \overleftarrow{\mathcal{D}}^\mathrm{F} \,.
\end{equation}
这对应一种``物理''的重正化方案，也就是说，算符
\begin{equation}
  \begin{split}
    \mathring{S}(t,x) = \zeta_\chi(t,\mu)\,S(t,x)\,,\qquad\mbox{其中}\quad
    \zeta_\chi(t,\mu) \equiv Z_\chi^{-1}\mathring{Z}_\chi(t,\mu)
    \label{eq:zetachi}
  \end{split}
\end{equation}
的反常量纲为零。\noeqn{eq:Zring} 中我们采用了记号
\begin{equation}
  \begin{split}
    \langle \mathcal{O}(t)\rangle \equiv \int\dd^4 x\,
    \langle \mathcal{O}(t,x)\rangle
    %% \label{eq:}
  \end{split}
\end{equation}
表示算符 $\mathcal{O}(t,x)$ 真空期望值的零动量傅里叶变换，下文沿用此约定。

上文引入的算符 $E(t,x)$、$S(t,x)$ 与 $R(t,x)$ 的费曼规则，按与
\sct{ssec:Feynman_rules} 所述相同的方式得到。把它们用流动场 $B(t,x)$ 与
$\chi_f(t,x)$ 表达：算符 $E(t,x)$ 给出一个双线性、一个三线性与一个四次胶子
顶点，如 \noeqn{fr:egg}--\noeqn{fr:egggg} 所示；$S(t,x)$ 对应 \noeqn{fr:sqq}
给出的单个双线性夸克顶点；$R(t,x)$ 对应 \noeqn{fr:rqq} 的双线性夸克顶点与
\noeqn{fr:rqqg} 的夸克-胶子顶点。与流顶点类似，连接到这些顶点的线用虚线箭头
画出，表示它们既可为流线也可为传播子。与流顶点不同的是，这些顶点中所有线的
流时间都指向相应算符。因此，这些顶点设定费曼图中最大的流时间。


利用 \noeqn{sec:pert} 节与 \ref{sec:implementation} 节所述的方法，我们可以
把 $E(t,x)$、$S(t,x)$ 与 $R(t,x)$ 的真空期望值计算到三圈。这些量的示例图见
\fig{fig:example_diagrams}。每种情形的图的数目列于表\,\ref{tab:numdias}；
但请注意，其中计入了 \sct{sec:implementation} 所述把四胶子顶点改写为三线性
顶点的做法。未计入其中的是含闭合流线圈的图，但含无标度子圈的积分被计入在
内。

\begin{figure}
  \begin{center}
    \begin{tabular}{ccc}
       \includegraphics[]{images/E_LO.pdf} (a) &
         \includegraphics[]{images/S_NLO_example.pdf} (b) &
         \includegraphics[]{images/E_NNLO_example.pdf} (c) \\[2em]
         \includegraphics[]{images/E_NNLO_example2.pdf}\hspace*{-1em}
         (d)\hspace*{1em} &
         \includegraphics[]{images/S_NNLO_example.pdf} (e) &
         \includegraphics[]{images/R_NNLO_example.pdf} (f) \\
    \end{tabular}
  \end{center}
	\caption{单圈处对 $\langle E(t)\rangle$ 有贡献的唯一图（a）；$\langle
    S(t)\rangle$ 与 $\langle R(t)\rangle$ 的一个两圈图（b）；以及四个三圈
    图：$\langle E(t)\rangle$ 的两个图（c,d）；$\langle S(t)\rangle$ 与
    $\langle R(t)\rangle$ 的一个图（e）；$\langle R(t)\rangle$ 的一个图
    （f）。}
	\label{fig:example_diagrams}
\end{figure}


\begin{table}
  \begin{center}
    \caption[]{\label{tab:numdias}本节所计算的量涉及的费曼图数目。三圈层次上，
      ``quenched \qcd''（标记为 $\TR=0$）的数目单独列出。}

    \begin{tabular}{|c||c|c|c|c|}
      \hline
      圈数 & 1 & 2 & 3 & 3  ($\TR=0$)\\\hline\hline
      $\langle E(t)\rangle$ & 1 & 11 & 232 & 211 \\\hline
      $\langle S(t)\rangle$ & 1 & 8 & 210 & 202 \\\hline
      $\langle R(t)\rangle$ & 1 & 11 & 311 & 300 \\\hline
    \end{tabular}
  \end{center}
\end{table}


\subsection{三圈胶子凝聚的结果}

我们把胶子凝聚的微扰结果写成
\begin{equation}
  \langle E(t) \rangle = \frac{3\alphas}{4\pi t^2} \frac{\NA}{8}
  \left[e_0 + \frac{\alphas}{4\pi} e_1 + \left(\frac{\alphas}{4\pi}\right)^2 e_2 +
    \mathcal{O}(\alphas^3)\right] + \mathcal{O}(m)\,,
  \label{eq:et}
\end{equation}
其中 $\alphas=\alphas(\mu)$，$\mu$ 为重正化标度。如图所示，此处忽略夸克质量
项；它们在 \nlo\ 的效应已在 \citere{Harlander:2016vzb} 中研究过，发现可以忽
略。此时，各 $e_i$ 仅依赖乘积 $z\equiv \mu^2t$。出于稍后即明的原因，用变量
\begin{equation}
  \begin{split}
    L(z) \equiv \ln(2z) + \EulerGamma\,,
    %% \label{eq:}
  \end{split}
\end{equation}
来参数化这一依赖较为方便。重正化标度的一个自然选择由梯度流的逆``涂抹半
径''（smearing radius）$q_8\equiv 1/\sqrt{8t}$ 给出。然而，显式的高阶计算表
明，稍大一点的取值
\begin{equation}
  \begin{split}
    \mu_0 = \frac{e^{-\EulerGamma/2}}{\sqrt{2t}} \approx 1.5\,q_8\,,
    \label{eq:mu0}
  \end{split}
\end{equation}
（对应 $L(\mu_0^2t)=0$）能改善微扰修正的稳定性\,\cite{Harlander:2018zpi}
（另见 \citere{Harlander:2016vzb}）。\citere{Iritani:2018idk} 利用
\citere{Harlander:2018zpi} 的两圈结果计算热力学量时其 $t\to 0$
外推的行为也印证了这一点。

于是我们把 \noeqn{eq:et} 的系数写成
\begin{equation}
  \begin{split}
    e_0(z) &= e_{0,0}\,,\qquad
    e_1(z) = e_{1,0} + \beta_0\,L(z)\,,\\
    e_2(z) &= e_{2,0} + (2\beta_0\,e_{1,0}+\beta_1)\,L(z) +
    \beta_0^2\,L^2(z)\,,
    \label{eq:efull}
  \end{split}
\end{equation}
其中 $\beta_0$、$\beta_1$ 取自 \noeqn{eq:anombg}。对各系数 $e_{i,j}$，我们
得到
\begin{equation}
  \begin{split}
    e_{0,0} =&\ 1\,,\qquad
    e_{1,0} = \left(\frac{52}{9} + \frac{22}{3} \ln 2 - 3 \ln 3\right)
    \CA - \frac{8}{9} \TF\\
    e_{2,0} =&\ 27.9786\, \CA^2  - (31.5652\ldots)\, \TF\CA
    + \left(16 \zeta(3) - \frac{43}{3}\right)
    \TF\CF + \left(\frac{8\pi^2}{27} -
    \frac{80}{81}\right) \TF^2\,,
    \label{eq:econst}
  \end{split}
\end{equation}
其中 $\zeta(z)$ 是黎曼 $\zeta$ 函数，$\zeta(3) =
1.20206\ldots$。$\TF\CA$ 系数中的三点号表明我们得到了解析形式的表达式，见附
录\,\ref{app:analytic} 的 \noeqn{eq:e20catf}。$\CA^2$ 系数的数值结果此处只
显示四位小数，而通过把单个积分的不确定度传播到最终结果所得到的数值精度估计
还要高出约六位。这一估计与分析结果已经算出的情形中数值与解析结果的观测之差
相符。下文列出的其他可观测量情况相同。\nlo\ 系数 $e_1$ 最先在
\citere{Luscher:2010iy} 中求得。取 $\mu=q_8$，可以发现 \nnlo\ 结果与
\citere{Harlander:2016vzb} 在亚百分比水平上一致。\footnote{准确地说，文献
  \citere{Harlander:2016vzb} 的全部色系数都与我们的新计算相符，唯独 $\CA^2$
  项``只有''前三位数字一致，这对应于该文献对该项过于乐观的不确定度估计。其
  对任何实际应用的影响应当无关紧要。}

对数项 $e_{n,k}L^k(z)$ 满足全阶递推公式（$1\leq k\leq n$）
\begin{equation}
  \begin{split}
    k\,e_{n,k} = \sum_{l=0}^{n-1}(n-l)e_{n-l-1,k-1}\beta_l\,,
  \end{split}
\end{equation}
它来自重正化群不变性，即
\begin{equation}
  \begin{split}
    \mu\dderiv{}{\mu}\langle E(t)\rangle = 0\,,\qquad
    \mu^2\dderiv{}{\mu^2}\alphas = \alphas\beta(\alphas) =
    -\frac{\alpha^2_s}{4\pi}\left(\beta_0 + \frac{\alphas}{4\pi}\beta_1
    + \ldots\right)\,.
    \label{eq:rgee}
  \end{split}
\end{equation}
这为我们的计算提供了一个受欢迎的检验。残留的 $\mu$ 依赖可用作微扰行为的探
针，此类研究可参见 \citere{Harlander:2016vzb}。

由于 $\langle E(t) \rangle$ 也与规范无关，在求值上述结果时我们可以自由选取
\qcd\ 规范参数的 Feynman 规范（$\xi=1$）与梯度流规范参数 $\kappa = 1$，这大
大便利了实际计算。不过我们也显式检验规范不变性：计算规范参数 $\xi$ 最高次幂
（即 $\xi^4$）的项并证明其系数为零。事实上，这在约化为主积分之后便已在代数
层面发生，甚至早于其数值代入。遗憾的是，保持规范参数完全一般会导致超出我可
们可用计算资源的中间表达式。

\subsection{三圈夸克凝聚的结果}

为得到加环场下的夸克凝聚，我们先考虑 \noeqn{eq:zetachi} 定义的从 \msbar\
重正化夸克场到加环夸克场的转换因子 $\zeta_\chi(t,\mu)$。运用
\sct{sec:pert} 所述的方法，我们可以对无质量夸克计算对格林函数 $\langle
R(t)\rangle$ 有贡献的图，最终得到结果
\begin{equation}
	\begin{split}
	  \zeta_\chi(t,\mu) = 1 &+ \frac{\alphas}{4\pi} \left(
          \frac{\gamma_{\chi, 0}}{2} L(\mu^2t) - 3 \CF \ln 3 - 4 \CF \ln
          2 \right) \\ &+ \left(\frac{\alphas}{4\pi}\right)^2 \Bigg\{
          \frac{\gamma_{\chi, 0}}{4} \left( \beta_0 +
          \frac{\gamma_{\chi, 0}}{2} \right) L^2 (\mu^2t) + \Big[
            \frac{\gamma_{\chi, 1}}{2} - \frac{\gamma_{\chi, 0}}{2}
            \left( \beta_0 + \frac{\gamma_{\chi, 0}}{2} \right) \ln 3
            \\ & \qquad \qquad - \frac{2}{3} \gamma_{\chi, 0} \left(
            \beta_0 + \frac{\gamma_{\chi, 0}}{2}\right) \ln 2 \Big]
          L(\mu^2t) + C_2 \Bigg\} + \order{\alphas^3,m} \,.
                \label{eq:zfchi}
	\end{split}
\end{equation}
\nlo\ 结果已在 \citere{Makino:2014taa} 中得到，依赖 $\mu$ 的 \nnlo\ 项则在
\citere{Harlander:2018zpi} 中导出。那篇文章还使用了系数 $C_2$
的一个初步结果，该系数在本文中首次导出。对各色系数保留四位小数，我们得到
\begin{equation}
  C_2 = -23.7947 \, \CA \CF + 30.3914 \, \CF^2 -(3.9226\ldots) \CF
  \TF\,,
  \label{eq:c2}
\end{equation}
其中同样如三点号所示，$\CF\TF$ 系数以解析形式获得，见附录\,\ref{app:analytic}
的 \noeqn{eq:c2cftf}。

现在转向夸克凝聚 $\langle S(t)\rangle$。由于它在手征极限 $m_f=0$ 下为零，
我们计算其按夸克质量 $m_f$ 展开的领头系数：
\begin{align}
  \left\langle \zeta_\chi S(t) \right\rangle \equiv \langle
  \mathring{S}(t) \rangle = \sum_{f=1}^{\NF} s_f(t) + \mathcal{O}(m^2)\,,\qquad
  s_f(t) \equiv m_f\dderiv{\langle
    \mathring{S}(t)\rangle}{m_f}\bigg|_{m=0}\,,
  \label{eq:sring}
\end{align}
其中理解为求导之后把\textit{所有}夸克质量置零。这一展开可以在圈积分与流时间
积分之前施行，也就是说，我们在夸克传播子中写出
\begin{equation}
  \begin{split}
    \frac{1}{\mathrm{i}\slashed{q}+m_f} =
    \frac{1}{\mathrm{i}\slashed{q}} + \frac{m_f}{q^2} +
    \mathcal{O}(m_f^2)\,,
    \label{eq:massexp}
  \end{split}
\end{equation}
并只保留被积函数中 $m$ 的领头线性项。所得积分仍是 \noeqn{eq:intrep} 的形
式。于是我们写
\begin{equation}
  s_f(t) = - \frac{\NC m_f}{8\pi^2t} \left[s_0 + \frac{\alphas}{4\pi} s_1 +
    \left(\frac{\alphas}{4\pi}\right)^2 s_2 +
    \mathcal{O}(\alphas^3)\right] + \mathcal{O}(m^2) \,.
\end{equation}
按 \noeqn{eq:renorm} 在 \msbar\ 方案中重正化夸克质量，我们得到
\begin{equation}
  \begin{split}
    s_0(z) &= s_{0,0}\,,\qquad
    s_1(z) = s_{1,0} + \frac{\gamma_{m,0}}{2}\,L(z)\,,\\
    s_2(z) &= s_{2,0} + \frac{1}{2}\left[(\gamma_{m,0}+2\beta_0)\,s_{1,0} +
      \gamma_{m,1}\right]L(z) +
    \frac{1}{8}(\gamma_{m,0}+2\beta_0)\gamma_{m,0}\,L^2(z)\,,
    \label{eq:sfull}
  \end{split}
\end{equation}
其中 $z=\mu^2t$，各 $\gamma_{m,i}$ 与 $\beta_i$ 取自 \noeqn{eq:anombg}。
非对数项为
\begin{equation}
  \begin{split}
    s_{0,0} &= 1\,,\qquad
    s_{1,0} = \left(4 + 4 \ln 2 - 3 \ln 3\right) \CF\\
    s_{2,0} &= 19.6422\, \CF^2 + 41.3897\, \CF \CA - (15.7975\dots)\, \TF \CF\,,
    \label{eq:sconst}
  \end{split}
\end{equation}
其中同样如原文所示，$\TF\CF$ 项解析已知（以三点号标示），可见附录
\,\ref{app:analytic} 的 \noeqn{eq:s20cftf}。\nlo\ 项已在
\citere{Makino:2018rys} 中得到，\nnlo\ 项是新结果。同样，此处只显示四位小
数，而我们对数值积分的不确定度估计还要高出若干位数字。

对数项 $s_{n,k}L^k(z)$ 满足全阶递推公式（$1\leq k\leq n$）
\begin{equation}
  \begin{split}
k\,s_{n,k} = \sum_{l=0}^{n-1}\left[\frac{\gamma_{n,l}}{2} +
  (n-l-1)\beta_l\right]s_{n-l-1,k-1}\,,
  \end{split}
\end{equation}
它来自重正化群不变性，即
\begin{equation}
  \begin{split}
    \mu\dderiv{}{\mu} s_f(t) = 0\,,\qquad
    \mu\dderiv{}{\mu}m_f = m_f\gamma_m(\alphas) =
    - m_f\frac{\alphas}{4\pi}\left(\gamma_{m,0} +
    \frac{\alphas}{4\pi}\gamma_{m,1} + \ldots\right)\,,
    %% \label{eq:}
  \end{split}
\end{equation}
以及 \noeqn{eq:rgee} 给出的 $\alphas$ 的重正化群方程。这再次为我们的计算提供
了一个受欢迎的检验。残留的 $\mu$ 依赖可用作微扰行为的探针。\fig{fig:S_plot}
展示了中心能标取 $\mu_0=(2te^{\EulerGamma})^{-1/2}$ 三个数值时的情形，参见
\noeqn{eq:mu0}。我们观察到与 \citere{Harlander:2016vzb} 中对 $\langle
t^2E(t)\rangle$ 发现的定性相似的行为：对较大的 $\mu_0$，随微扰阶增高重正化
标度依赖的减小相当显著；而当 $\mu_0$ 更接近非微扰区域时，这一改善有所减弱。
若取 $\mu$ 在中心标度附近变化因子二作为缺失更高阶导致的不确定度估计，则对全
部三个 $\mu_0$ 值，相继微扰阶的区间彼此重叠。然而在这种情形下，曲线的趋势并
未明确支持以 $\mu_0$ 作为中心标度，而是似乎偏向稍小的标度。

\begin{figure}
  \begin{center}
    \begin{tabular}{cc}
      \includegraphics[width=.45\textwidth]{images/Sring_3.pdf} &
      \includegraphics[width=.45\textwidth]{images/Sring_10.pdf}\\
      \hspace*{2em}(a) & \hspace*{2em}(b)\\
      \multicolumn{2}{c}{\includegraphics[width=.5\textwidth]{%
          images/Sring_130.pdf}}\\
      \multicolumn{2}{c}{\hspace*{2em}(c)}
    \end{tabular}
  \end{center}
	\caption{$s_f(t)$（定义见 \noeqn{eq:sring}）的重正化标度依赖，中心标度
          $\mu_0=e^{-\EulerGamma/2}/\sqrt{2t}$ 取三个不同数值。
          (a)~$\mu_0=3$\,GeV，对应 $t=(0.03\,\mathrm{fm})^2$，取 $\NF = 3$
          与 $\alpha_\mathrm{s}^{(3)}(3\,\mathrm{GeV}) =
          0.248$；
          (b)~$\mu_0=10$\,GeV（$t=(0.01\,\mathrm{fm})^2$），取 $\NF = 5$ 与
          $\alpha_\mathrm{s}^{(5)}(10\,\mathrm{GeV}) = 0.178$；
          (c)~$\mu_0=130$\,GeV（$t=(8\cdot
          10^{-4}\,\mathrm{fm})^2$），取 $\NF = 5$ 与
          $\alpha_\mathrm{s}^{(5)}(130\,\mathrm{GeV}) = 0.112$。
          $\alphas(\mu)$ 与 $m_f(\mu)$ 的跑动借助
          \texttt{RunDec}\,\cite{Chetyrkin:2000yt,Herren:2017osy}
          在相应圈数下求值。所有曲线均归一到 $\mu=\mu_0$ 处的 \nnlo\
          结果。}
	\label{fig:S_plot}
\end{figure}

另一个重要的自洽性检验是：$\langle S(t) \rangle$ 与
$\langle R(t) \rangle$ 可以用 \citere{Harlander:2018zpi} 中得到的同一个流动
夸克场重正化常数 $Z_\chi$ 进行重正化。$\langle S(t) \rangle$ 与 $\langle
R(t)\rangle$ 也都与规范无关。为计算完整的 \nnlo\ 结果，我们同样选取 Feynman
规范（$\xi=1$）与 $\kappa=1$。我们显式检验规范不变性：计算规范参数 $\xi$
最高次幂的项——对 $\langle S(t) \rangle$ 与 $\langle R(t)\rangle$ 为
$\xi^3$——并证明其系数为零；这在约化为主积分之后立即发生。
